% ------------------------------------------------------------------
% Chapter: Conclusion
% ------------------------------------------------------------------
\chapter{Conclusion}

Diese Arbeit verbindet klassische Bilanzlogik mit moderner, reproduzierbarer
Datenanalyse zu einem kompakten Frühwarnsystem. Kern ist eine Pipeline, die
(a) Multikollinearität systematisch abbaut (VIF), (b) Features über Filter-,
Wrapper- und Embedded-Methoden selektiert und deren Robustheit mit Nested-CV
und dem Stabilitätsmaß nach \textcite{nogueira2018stability} prüft, (c) EPV-
Leitplanken berücksichtigt und (d) konsensorientierte Endmengen bildet. Auf
Basis dieser Sets zeigt der Modellvergleich über fünf Horizonte, dass ein
Random Forest in vier von fünf Fällen die höchste ROC-AUC erzielt, während ein
Soft-Voting-Ensemble bei kurzer Frist (H1) knapp vorn liegt. Bemerkenswert ist,
dass kleine, stabile Merkmalsmengen die Güte der Vollmengen weitgehend halten
(AUC-Retention ca. 97--102\%) und damit Interpretierbarkeit und
Verallgemeinerbarkeit fördern.

Methodisch tragen drei Entscheidungen wesentlich zur Glaubwürdigkeit bei:
(1) klare Trennung von Auswahl und Bewertung (Nested-CV in der Selektion,
stratifizierte CV in der Modellierung), (2) Stabilität als Qualitätskriterium
für Merkmalslisten, (3) EPV-Disziplin, um Überparametrisierung zu vermeiden.
Diese Maßnahmen reduzieren Optimismus und erleichtern eine ehrliche
Projektdokumentation.

Praktisch bietet die Pipeline einen Einstieg für Kredit- und
Risikofunktionen, die auf interpretierbare Kennzahlen, schlanke Modelle und
stabile Entscheidungen angewiesen sind. Künftige Arbeit (Kap.~\ref{chap:future-work})
liegt in gruppenbewusster CV (sofern Firmen-IDs verfügbar),
Kalibrierung und Schwellenwahl mit Kostenfunktion, systematischem
Hyperparameter-Tuning, stabilitätsgewichteter Konsensbildung sowie temporalen
Validierungsschemata. Damit lässt sich die hier gezeigte Robustheit weiter in
Richtung operativer Umsetzung entwickeln.
